% 致谢
\chapter*{致\hspace{1em}谢}
\addcontentsline{toc}{chapter}{致谢}
\markboth{致谢}{致谢}

本论文的完成离不开众多人的帮助与支持。

首先，感谢我的指导教师金蒙豪在我本科学习期间对我的认真指导，从大一开始的SRTP的导师课题互选开启了我的本科的科研经历，一步步教导我从一个对科研一片茫然的本科生变成一个有一定经验和科研能力的本科毕业生。最开始，金老师带我引入了关于微磁学的研究，初步研究了有关Skyrmions的一些内容，并且教导我能够利用OOMMF来进行最简单的微磁学仿真，并指导了我的一篇基于Skyrmions的人工神经元的构建的专利的内容。在此之后，又给我开拓了"三维Skyrmions"——Hopfion的研究，才让我接触到这个科学界最前沿的尖端领域的知识，也更进一步学会了Mumax3这种可视化更好的微磁学仿真软件，让我在微磁学领域有了更多的知识储备和了解，对我之后的研究打下了无比扎实的基础，最终影响到我的最后本科毕业设计的研究。关于本次的本科毕业设计，从课题的选定、研究方案的设计到论文的撰写修改，金老师都给予了耐心的指导和很好的建议，让我在磁性拓扑结构这一前沿领域获得了很多研究经验。

其次，感谢杭州电子科技大学电子信息学院第三教学楼519实验室为本课题提供的计算资源和科研条件，在最开始我缺乏算力资源的阶段，给我提供了微磁学仿真最基础的条件和算力，才能让我完成最终的毕业设计的产出。感谢本科期间课题组全体成员在日常学习和研究中的交流与帮助，大家的支持对我帮助很大。

最后，特别感谢所有在本科四年学习生活中给予我帮助的老师和同学们和我的家人，感谢他们一直以来的理解、支持与鼓励。
